\chapter{Ngô Bào Châu}\label{ch2010ngobaochau:nguxf4-buxe0o-chuxe2u-fields-medal-2010-a-comprehensive-report-on-his-work-and-impact}

\begin{center}\rule{0.5\linewidth}{0.5pt}\end{center}

\textbf{Prepared:} July 2026 \textbf{Subject:} Ngô Bào Châu (b. June 28, 1972, Hanoi, Vietnam) \textbf{Position:} Professor, Department of Mathematics, University of Chicago; Scientific Director, Vietnam Institute for Advanced Study in Mathematics \textbf{Doctoral Advisor:} Gérard Laumon \textbf{Fields Medal Citation:} \emph{``For his proof of the Fundamental Lemma in the theory of automorphic forms through the introduction of new algebro-geometric methods.''}

\begin{center}\rule{0.5\linewidth}{0.5pt}\end{center}

\begin{prereqs} Algebraic geometry (stacks, intersection theory, perverse sheaves), representation theory (reductive groups, Hecke algebras), number theory (automorphic forms, trace formula). \end{prereqs}

\begin{objectives} State the Fundamental Lemma \cite{ngo2010} and explain its role in the trace formula; describe the geometric interpretation via affine Springer fibers; understand the role of the endoscopic classification. \end{objectives}

\section{Problem History}\label{ch2010ngobaochau:problem-history}

The \textbf{Langlands program} is a vast web of conjectures connecting number theory, representation theory, and algebraic geometry. A central pillar is the \textbf{trace formula}, introduced by Arthur and Selberg, which relates the spectral decomposition of automorphic forms to the orbital integrals of conjugacy classes in reductive groups. The trace formula is a powerful tool for establishing functoriality and for proving reciprocity laws.

The \textbf{Fundamental Lemma} is a combinatorial identity that arises when comparing the trace formula for a group $G$ with the trace formula for an endoscopic group $H$ of $G$. Precisely, for a reductive group $G$ over a local field $F$ and an endoscopic group $H$, the Fundamental Lemma asserts the identity of certain orbital integrals on $G$ with related orbital integrals on $H$, under the local Langlands correspondence.

The conjecture was formulated by Langlands and Shelstad in 1987 in the context of \emph{stabilization} of the trace formula. For $G = \text{GL}_n$, the lemma was proved by Kottwitz (1988). For general reductive groups, the conjecture was open for over 20 years, standing as a major technical obstacle to the stabilization program.

The endoscopic groups of $G$ are defined through the notion of $L$-parameters. An endoscopic group $H$ of $G$ comes with a Langlands dual group ${}^L H$ and a natural map ${}^L H \to {}^L G$ that is injective on Weyl groups. The Fundamental Lemma compares orbital integrals on $G$ with orbital integrals on $H$ that are ``transferred'' using test functions related through the local Langlands correspondence. For split groups like $\text{GL}_n$, the endoscopic groups are products of general linear groups, but for non-split groups, the endoscopic structure is considerably more complicated.

\section{How Was It Solved}\label{ch2010ngobaochau:how-was-it-solved}

\textbf{The Geometric Approach.}

Ngô Bào Châu proved the Fundamental Lemma by introducing a profound \emph{geometric interpretation} of the orbital integrals. The key idea is to view the orbital integrals as \textbf{intersection numbers} on certain algebraic varieties called \textbf{affine Springer fibers}.

For a reductive group $G$ over a local field $F$ with a uniformizer $\varpi$, the \textbf{affine Springer fiber} $\mathcal{X}_\gamma$ is the variety

\[
\mathcal{X}_\gamma = \{ g \in G(\mathbb{C}[[t]]) / G(\mathbb{C}[[t]]) : g^{-1} \gamma g \equiv \gamma \mod t^n \},
\]

where $\gamma \in G(F)$ is a semisimple element and $t$ is the uniformizer. The orbital integral of a test function $f$ with respect to $\gamma$ is proportional to the Euler characteristic (with compact support) of the corresponding affine Springer fiber:

\[
\text{Orb}_\gamma(f) \propto \chi_c(\mathcal{X}_\gamma, \mathcal{L}_f),
\]

where $\mathcal{L}_f$ is a certain perverse sheaf (or local system) determined by $f$.

The affine Springer fiber $\mathcal{X}_\gamma$ sits inside the affine Grassmannian $\operatorname{Gr}_G = G(\mathbb{C}((t))) / G(\mathbb{C}[[t]])$, which is an ind-variety that plays a central role in geometric representation theory. The key property is that $\mathcal{X}_\gamma$ has finite-dimensional cohomology, and its Euler characteristic encodes the orbital integral.

\textbf{The Trace Formula as Intersection Theory.}

Ngô showed that the Fundamental Lemma is equivalent to a statement about the \emph{intersection multiplicities} of cycles on the affine Springer fiber. The comparison between the orbital integrals for $G$ and $H$ becomes a comparison between intersection numbers on $\mathcal{X}_\gamma^G$ and $\mathcal{X}_\gamma^H$, which Ngô established using:

\begin{enumerate}
  \item A \textbf{fibration formula} for the Euler characteristic of affine Springer fibers, expressing the Euler characteristic as a sum over fixed points of a torus action.
  \item A \textbf{geometric lemma} relating the fixed-point sets on $\mathcal{X}_\gamma^G$ to those on $\mathcal{X}_\gamma^H$, using the structure of the endoscopic group $H$ and the \textbf{Springer resolution}.
  \item The \textbf{localization theorem} in equivariant cohomology, which reduces the computation of intersection numbers to a sum over isolated fixed points.
\end{enumerate}

The crucial insight is that the Fundamental Lemma, which is a purely combinatorial identity at the level of orbital integrals, has a \emph{geometric origin}: it arises from the decomposition of intersection numbers on affine Springer fibers under the action of the endoscopic group.

\textbf{Endoscopic Classification.}

The endoscopic classification, due to Arthur, describes the automorphic representations of $G$ in terms of those of the endoscopic groups $H$. The Fundamental Lemma is needed to compare the trace formulas for $G$ and $H$, and hence to identify the relevant automorphic representations. Ngô's proof establishes this comparison for all reductive groups, completing a key step in Arthur's program.

Arthur's stabilization of the trace formula expresses the trace formula for $G$ as a sum of stabilized trace formulas for endoscopic groups $H$, and the Fundamental Lemma is the key identity that makes this decomposition work. With Ngô's proof in hand, the stabilization is established, and the full force of the trace formula can be applied to the Langlands functoriality problem.

\section{Impact of This Discovery}\label{ch2010ngobaochau:impact-of-this-discovery}

Ngô's proof of the Fundamental Lemma has had a profound impact on the Langlands program. It immediately implies the \textbf{stabilization} of the trace formula for all reductive groups, which in turn yields:

\begin{enumerate}
  \item The \textbf{Langlands--Shahidi method} for establishing functoriality for certain classes of representations.
  \item The \textbf{global Gan--Gross--Prasad conjecture} for unitary groups, connecting special values of $L$-functions to periods of automorphic forms.
  \item The \textbf{base change} and \textbf{automorphic induction} for classical groups.
\end{enumerate}

The geometric methods introduced by Ngô---particularly the use of affine Springer fibers and their intersection theory---have opened new directions in geometric representation theory. The theory of \emph{exotic} sheaves on affine Grassmannians, developed by Lusztig and Ngô, provides new tools for studying the local Langlands correspondence.

The work also connects to \emph{geometric Langlands} and the study of perverse sheaves on moduli spaces of $G$-bundles, suggesting deep connections between the arithmetic Langlands program and the geometric Langlands program. The affine Springer fibers serve as a bridge between the two, and their intersection theory encodes information about both the arithmetic and geometric sides.

The \emph{orbital integral} of a test function $f$ on $G(F)$ with respect to a semisimple element $\gamma$ is defined as:
\[
\operatorname{Orb}_\gamma(f) = \int_{G(F)/T_\gamma(F)} f(g^{-1} \gamma g) \, dg,
\]
where $T_\gamma$ is the centralizer of $\gamma$ in $G$ and $dg$ is the Haar measure on $G(F)$ normalized so that $G(\mathcal{O}_F)$ has volume 1. The Fundamental Lemma asserts that for test functions $f$ on $G(F)$ and $f_H$ on $H(F)$ related by the local Langlands correspondence, the orbital integrals satisfy:
\[
\operatorname{Orb}_\gamma(f) = \sum_{\gamma' \in \operatorname{Emb}(\gamma, H)} c(\gamma, \gamma') \operatorname{Orb}_{\gamma'}(f_H),
\]
where the sum is over embeddings of $\gamma$ into $H$ and $c(\gamma, \gamma')$ are explicit constants depending only on the local Langlands parameters. Ngô's proof establishes this identity by interpreting both sides as intersection numbers on affine Springer fibers.

The global trace formula, which compares the spectral decomposition on $G(\mathbb{A})/\Gamma$ with the orbital integrals of conjugacy classes, requires a global version of the Fundamental Lemma. This ``global Fundamental Lemma'' compares global orbital integrals---integrals over $G(\mathbb{A})$ with respect to the Tamagawa measure---and Ngô's local methods provide the starting point for this global comparison. The global version involves additional complications from the archimedean places and from the interaction between different local factors.

\section{Modern Views and Improvements}\label{ch2010ngobaochau:modern-views-and-improvements}

The Fundamental Lemma has been generalized in several directions. The \textbf{relative trace formula} (RTF), which compares orbital integrals for $G$ relative to a subgroup $H$, requires a relative version of the Fundamental Lemma. This has been established in several cases by Ngô and others, with applications to the arithmetic Gan--Gross--Prasad conjecture and the study of periods of automorphic forms.

The \emph{affine Grassmannian} $\operatorname{Gr}_G = G(\mathbb{C}((t))) / G(\mathbb{C}[[t]])$ is an ind-variety that parameterizes lattices in $\mathbb{C}((t))^n$ (for $G = \operatorname{GL}_n$) with additional structure. The affine Springer fiber $\mathcal{X}_\gamma$ is a subvariety of $\operatorname{Gr}_G$ defined by the condition that the corresponding lattice is stable (up to modification) under conjugation by the semisimple element $\gamma$. The key property of $\mathcal{X}_\gamma$ is that it has finite-dimensional cohomology, and its Euler characteristic encodes the orbital integral of $\gamma$.

The \emph{intersection pairing} on the affine Grassmannian is defined using the convolution product: for two perverse sheaves $\mathcal{F}$ and $\mathcal{G}$ on $\operatorname{Gr}_G$, the convolution $\mathcal{F} * \mathcal{G}$ is a perverse sheaf on $\operatorname{Gr}_G$ whose stalk at a point $g$ is given by the hypercohomology of the convolution product over the intersection $\operatorname{Gr}_G \times_{\operatorname{Gr}_G} \operatorname{Gr}_G$. The Fundamental Lemma can be reformulated as a statement about the convolution product: it asserts that the convolution of the characteristic functions of certain double cosets on $G$ and $H$ are compatible under the local Langlands correspondence. Ngô's geometric proof translates this compatibility into a statement about intersection numbers on affine Springer fibers.

Open problems include: proving the relative Fundamental Lemma in full generality, establishing the \textbf{global} Fundamental Lemma (which compares global orbital integrals), and extending the geometric methods to \emph{non-split} reductive groups and \emph{mixed-characteristic} fields.

The \emph{relative Fundamental Lemma} compares relative orbital integrals for a pair $(G, H)$ where $H$ is a subgroup of $G$, rather than comparing orbital integrals on $G$ with those on an endoscopic group. The relative trace formula, which arises from this comparison, has applications to the arithmetic Gan--Gross--Prasad conjecture (relating special values of $L$-functions to periods of automorphic forms) and to the study of equidistribution of rational points on Shimura varieties. Ngô's geometric methods have been extended to the relative setting by several authors, but the full generality of the relative Fundamental Lemma remains open.

The \emph{mixed-characteristic} setting, where the base field has characteristic 0 but residue characteristic $p > 0$, presents additional challenges for Ngô's methods. The affine Grassmannian and affine Springer fibers have a different structure in the mixed-characteristic case, and the intersection theory must be replaced by $p$-adic Hodge theory. Recent work by Fargues and Scholze on the Fargues--Fontaine curve provides a geometric framework for the mixed-characteristic Langlands correspondence, and Ngô's methods may find applications in this setting.

Ngô continues to work on the \emph{geometric foundations} of the Langlands program, developing the theory of perverse sheaves on affine Grassmannians and their role in the local and global Langlands correspondence.

The \emph{exotic sheaves} on the affine Grassmannian, introduced by Lusztig and Ngô, are perverse sheaves that encode the representations of the Langlands dual group. The exotic sheaves are defined using the convolution product on the affine Grassmannian, and their endomorphism algebras are isomorphic to the universal enveloping algebra of the Langlands dual group. This provides a geometric realization of the representation theory of reductive groups, and it connects Ngô's work on the Fundamental Lemma to the geometric Satake correspondence.

The \emph{parahoric Hecke algebra} is a variant of the usual Hecke algebra that plays a role in the local Langlands correspondence for groups with parahoric subgroups. Ngô's geometric methods have been applied to compute the structure of parahoric Hecke algebras, with applications to the study of automorphic representations on quaternion algebras and other non-split groups.

Recent work by Fargues and Scholze on the Fargues--Fontaine curve and the local Langlands correspondence suggests a new geometric framework where Ngô's methods may find further applications. The relationship between affine Springer fibers and the geometry of the Fargues--Fontaine curve is an active area of research.

The \emph{geometric Satake correspondence} relates the category of perverse sheaves on the affine Grassmannian $\operatorname{Gr}_G$ to the representation theory of the Langlands dual group ${}^L G$. More precisely, the equivalence $\text{Perv}(\operatorname{Gr}_G) \cong \text{Rep}({}^L G)$ sends the intersection cohomology complex $\operatorname{IC}_{\overline{\operatorname{Gr}_G^\lambda}}$ to the irreducible representation $V_\lambda$ of ${}^L G$ with highest weight $\lambda$. The convolution product on $\operatorname{Gr}_G$ corresponds to the tensor product of representations, and the stalks of the IC complexes encode the multiplicities of weight spaces. Ngô's work on the Fundamental Lemma uses the geometric Satake correspondence to translate the combinatorial identity into a statement about intersection numbers on affine Springer fibers.

The \emph{Springer fiber} $\mathcal{B}_\gamma$ for a semisimple element $\gamma \in G$ is the variety of Borel subgroups containing $\gamma$. The Springer fiber has a natural action of the centralizer $T_\gamma$ of $\gamma$, and the cohomology $H^*(\mathcal{B}_\gamma)$ carries a representation of the Weyl group $W$ of $G$ (the \emph{Springer representation}). The affine Springer fiber $\mathcal{X}_\gamma$ is the affine analogue: it parameterizes lattices (in the $\operatorname{GL}_n$ case) that are stable under $\gamma$ up to modification. The key insight of Ngô's proof is that the orbital integral $\operatorname{Orb}_\gamma(f)$ can be expressed as an intersection number on $\mathcal{X}_\gamma$, and the comparison between $G$-orbital integrals and $H$-orbital integrals becomes a comparison of intersection numbers on $\mathcal{X}_\gamma^G$ and $\mathcal{X}_\gamma^H$.

\begin{center}\rule{0.5\linewidth}{0.5pt}\end{center}

\section*{Summary}
\addcontentsline{toc}{section}{Summary}

Ngô Bào Châu proved the Fundamental Lemma of the Langlands program by establishing a geometric interpretation of orbital integrals as intersection numbers on affine Springer fibers. The proof uses equivariant intersection theory and the structure of endoscopic groups to reduce the combinatorial identity to a comparison of Euler characteristics. The Fundamental Lemma is a key ingredient in the stabilization of the trace formula, with applications to functoriality, the Gan--Gross--Prasad conjecture, and the study of automorphic $L$-functions.

\begin{center}\rule{0.5\linewidth}{0.5pt}\end{center}
